\chapter*{\centering EXECUTIVE SUMMARY}
\addcontentsline{toc}{chapter}{EXECUTIVE SUMMARY}

Health insurance penetration in Cambodia is very low, and the insurers that do operate there underwrite new applicants with static, rule-based systems largely imported from mature markets. These rules are fixed at deployment: they ignore interactions between risk factors, cannot adapt as the applicant pool shifts, and offer no built-in safeguard against demographic bias. This thesis asks whether \emph{contextual bandits} --- a family of online reinforcement-learning algorithms that learn from the outcome of every decision they make --- can replace those static rules, and whether lightweight statistical guardrails can keep the resulting portfolio fair and actuarially sound.

The project designs, implements, and rigorously evaluates a complete adaptive-underwriting pipeline on a synthetic dataset of 2{,}000 Cambodian applicants calibrated to the Cambodia Demographic and Health Survey. Four pre-registered experiments --- 20 independent random seeds, 5{,}000 decisions each, with paired significance testing --- compare the learners against a complete ladder of baselines and stress-test their fairness, robustness, and operational cost.

\bigskip
\noindent\textbf{Key Contributions}
\begin{enumerate}
    \item \textbf{An adaptive underwriting engine.} Three contextual-bandit algorithms (LinUCB, LinTS, and Epsilon-Greedy) decide, for each applicant, among four actions --- \emph{standard}, \emph{rated}, \emph{decline}, and \emph{refer} --- and update their model after observing the actuarial profit of each decision.
    \item \textbf{A Cambodia-calibrated test environment.} A profit-based actuarial reward simulator and a 2{,}000-applicant synthetic dataset, anchored on national demographic and disease-prevalence statistics, provide a reproducible benchmark in which the optimal decision for every applicant is known.
    \item \textbf{Built-in fairness monitoring.} Sliding-window Population Stability Index (PSI) guardrails and the U.S.\ EEOC ``four-fifths'' rule jointly audit the approved portfolio for regional and occupational disparity at every step of the run.
    \item \textbf{A human-in-the-loop safety layer.} A wrapper escalates only the genuinely ambiguous cases to a human underwriter, capturing expert judgement while keeping manual-review cost low.
\end{enumerate}

\bigskip
\noindent\textbf{Principal Findings}
\begin{itemize}
    \item \textbf{Bandits beat every deployable alternative.} LinUCB earns \$90{,}540 in cumulative reward versus \$72{,}292 for the static XGBoost rule --- a \textbf{25\%} improvement (Wilcoxon $p < 0.001$, Cohen's $d = 2.98$, large effect) --- and LinTS improves on the rule by roughly 30\%. Late-run regret falls to \$2.20 per decision against the rule's \$9.01.
    \item \textbf{Adaptation does not cost fairness.} With no explicit fairness term in its objective, the bandit holds regional and occupational approval parity at 85.7\% and 90.1\% (both clearing the 80\% four-fifths threshold) and keeps the maximum sliding-window PSI in the GREEN/low-AMBER band throughout the run.
    \item \textbf{The learning that matters is online estimation, not exploration.} An ablation shows that continuously re-fitting per-action coefficients --- rather than the exploration bonus --- is the load-bearing mechanism, and that the adaptive learners absorb a simulated demographic shock far better than the static rule.
    \item \textbf{Human oversight pays for itself.} Routing only the hardest 1.5\% of cases to a human underwriter raises cumulative reward by a further 6.5\%, at a review cost of just 2.6\% of reward.
    \item \textbf{An honest ceiling.} A post-hoc analysis found that a trivial constant policy --- rating \emph{every} applicant at a flat surcharge --- scores higher still, but is neither commercially viable nor regulatorily admissible. It is reported as a ceiling that \emph{bounds}, but does not overturn, the contribution: the bandits beat every alternative an insurer could actually deploy.
\end{itemize}

\bigskip
\noindent\textbf{Significance and Impact}

\noindent Traditional emerging-market underwriting applies rules calibrated on foreign loss experience to populations whose disease profiles and occupational structures differ fundamentally. A contextual bandit instead learns directly from the population it serves and updates continuously as that population evolves --- while remaining light enough to run on the low-resource mobile infrastructure typical of Cambodian insurers. Combined with PSI guardrails and human escalation for borderline cases, the architecture offers a path to underwriting that is at once more accurate, more fair, and more adaptable than the static systems it would replace.

This work is a proof of concept on synthetic data: its value lies not in the specific figures --- which will differ on real claims data --- but in demonstrating that the full pipeline (data generation, bandit training, fairness auditing, benchmark comparison, human integration, and drift monitoring) is feasible within the resource and regulatory constraints of an emerging-market insurer. The natural next step is a controlled pilot with an industry partner.
